\documentclass[11pt]{article}
\usepackage[a4paper,margin=1in]{geometry}
\usepackage[T1]{fontenc}
\usepackage{lmodern,amsmath,amssymb,amsthm,booktabs,microtype}
\usepackage[hidelinks]{hyperref}
\hypersetup{pdftitle={Sparse Initial Conditions and the Trace Space of Langton's Ant},pdfauthor={Hao Ke}}
\setlength{\emergencystretch}{2em}
\newtheorem{theorem}{Theorem}[section]
\newtheorem{lemma}[theorem]{Lemma}
\newtheorem{proposition}[theorem]{Proposition}
\newtheorem{corollary}[theorem]{Corollary}
\theoremstyle{definition}\newtheorem{definition}[theorem]{Definition}
\theoremstyle{remark}\newtheorem{remark}[theorem]{Remark}
\newcommand{\Tr}{\mathcal T}
\newcommand{\Pos}{\mathcal P}
\newcommand{\N}{\mathbb N}
\newcommand{\Z}{\mathbb Z}
\newcommand{\Rr}{\mathrm R}
\newcommand{\Ll}{\mathrm L}
\title{Sparse Initial Conditions and the Trace Space of Langton's Ant}
\author{Hao Ke}
\date{}
\begin{document}
\maketitle
\begin{abstract}
How many different observations can an ant with at most $m$ initially black
cells produce in its first $n$ updates, when the cell coordinates are unrestricted?
We prove that the number of anchored turn prefixes is $\Theta_m(n^m)$ for every
fixed $m$. The upper bound encodes a history by the times of its first black
reads. The matching lower exponent is realized by a protected single-cell
scattering channel, repeated with complete history retained. Every input in
this lower-bound family ultimately follows the standard highway.
For infinite streams, the countable compact trace spaces obey the exact
derived-set identity $\Tr_m'=\Tr_{m-1}$: their Cantor--Bendixson rank is $m+1$.
An explicit compact subfamily of eventually standard streams has the same
rank and prefix exponent. We also determine the optimal passive prediction
mistake budget to be $m$. These results classify sparse-source observation complexity uniformly over
unrestricted initial coordinates, transient times and future geometries.
\end{abstract}

\section{The anchored sparse-source question}
A fixed number of black cells can be placed arbitrarily far apart, so their
effect on an ant may first appear after a long flight. We ask how this sparse
initial information limits the number of different observations, without
restricting those distances or the subsequent trajectory. The object is a
family of traces from a common starting pose, rather than one selected orbit.

We use $\N=\{0,1,2,\ldots\}$.
A classical ant state is a complete finite black set $B\subset\Z^2$, a head
position $p$ and a heading $h\in\{N,E,S,W\}$. At each update a white read turns
right, a black read turns left, the current cell toggles, and the ant moves one
edge. We fix the initial pose to $(0,0),N$; translations and quarter turns
normalize all other poses without changing the turn word. Time $t$ is before
update $t$.

Let $\Tr_m\subset\{\Rr,\Ll\}^{\N}$ be the infinite turn streams of complete
inputs with at most $m$ black cells. Initial coordinates have no size bound.
Let $p_m(n)$ count their length-$n$ prefixes. Let $p_m^H(n)$ count prefixes from
those inputs whose future is eventually the standard period-104 highway.
A statement about $p_m$ describes the entire sparse-source family, rather than
the factor complexity of one selected trace.

\begin{theorem}[Sharp prefix exponent]\label{thm:prefix}
For all nonnegative integers $n,m$,
\begin{equation}\label{eq:upper}
 p_m^H(n)\le p_m(n)\le\sum_{j=0}^{\min(m,n)}\binom nj.
\end{equation}
For every $L\ge 0$ such that $11017+208L+20\min(m,L)\le n$,
\begin{equation}\label{eq:lower}
 \sum_{j=0}^{\min(m,L)}\binom Lj\le p_m^H(n).
\end{equation}
Consequently, for each fixed $m$,
\[
 p_m(n)=\Theta_m(n^m),\qquad p_m^H(n)=\Theta_m(n^m),\qquad
 \lim_{n\to\infty}\frac{\log p_m(n)}{\log n}=m.
\]
\end{theorem}

The exponent is exact; an exact leading coefficient is not claimed.
The uniform bounds \eqref{eq:upper}--\eqref{eq:lower} also permit $m$ to vary.
For $m=0$ there is just the blank-orbit prefix, so both counts equal one.

Equip $\{\Rr,\Ll\}^{\N}$ with its product (Cantor) topology: sets
specifying a finite prefix form a basis of sets that are both open and closed.
For a subspace $X$, let $X'$
be its set of accumulation points, and define $X^{(0)}=X$,
$X^{(j+1)}=(X^{(j)})'$. The rank convention below is the first empty derivative.

\begin{theorem}[Exact derivative hierarchy]\label{thm:topology}
Each $\Tr_m$ is countable and compact, and
\[
 \Tr_m'=\Tr_{m-1}\quad(m\ge 1),\qquad
 \Tr_m^{(m)}=\Tr_0,\qquad \Tr_m^{(m+1)}=\varnothing.
\]
Thus its Cantor--Bendixson rank is exactly $m+1$.
\end{theorem}

\begin{theorem}[Complexity inside the positive basin]\label{thm:positive}
There are explicit compact spaces $\Pos_m\subset\Tr_m$, all of whose streams
are eventually standard, such that $\Pos_m'=\Pos_{m-1}$ for $m\ge 1$,
$\Pos_0=\Tr_0$, and their prefixes satisfy \eqref{eq:lower}.
\end{theorem}

The last assertion separates observation complexity from convergence:
eventually standard inputs already achieve the largest polynomial exponent
and derived-set depth allowed by their initial mass. The spaces $\Tr_m$ are
not assumed shift-invariant. In particular we do not identify their anchored
prefix exponent with entropy of a bounded-mass dynamical system.

\paragraph{Prior methods and attribution.}
Gajardo~\cite{gajardo} introduced the symbolic read/turn projection for arbitrary
colorings and studied its subshift. Same-cell alternation and the corresponding
minimal first-read source are established semantics. The compactness and
encoding arguments below are elementary persistent-query arguments; the
particular dynamics enter the lower bound and the availability of arbitrarily
late fresh cells. The 228-update, phase-13 single-obstacle response used here
is already present in the one-black project~\cite{oneblack}. We reconstruct its
complete half-plane form rather than claim a new scattering response or a new
one-black theorem. The question addressed here is the sharp all-$m$ census and
hierarchy of anchored sparse-source traces, including simultaneous saturation
by positive inputs.

Cantor--Bendixson rank is also used to study countable sofic shifts, for
example by Salo and T\"orm\"a~\cite{salo}; our spaces are anchored and need
not be shift-invariant. Worst-case mistake counts are a classical criterion
for online prediction~\cite{littlestone}. The exact hierarchy and budget here
are proved directly from sparse first reads and the protected ant channel.

\section{First-read times encode the entire finite history}
For a prescribed word $w$, its positions and headings are obtained by turning
and moving without yet assigning colors. Consecutive reads of the same cell
must alternate between $\Rr$ and $\Ll$. This condition is sufficient: blacken
exactly the cells whose first read is $\Ll$, and start every other cell white.
An induction on time verifies the whole word.

Write $M(w)$ for this unique minimal source, and write $E(w)$ for the times of
its first black reads. Then $|E(w)|=|M(w)|$. Any other realizing source agrees
with $M(w)$ on the complete read footprint, but can contain additional unread
black cells. We retain those cells in physical execution.

A finite word is called feasible if it is realized from the fixed initial
pose by some finite initial black field. Equivalently, successive visits to
each cell request alternating colors.

\begin{lemma}[Time-code injection]\label{lem:code}
A feasible length-$n$ word is uniquely determined by $E(w)\subset\{0,\ldots,n-1\}$.
\end{lemma}
\begin{proof}
Starting from the normalized pose, at an unvisited cell choose $\Ll$ exactly
when the present time is in $E(w)$, otherwise choose $\Rr$. At a visited cell
use the opposite of its preceding read. Toggle and move. Each decision and
hence every future position is determined. A proposed code marking a revisit
is rejected. For a code from an actual word this decoder reproduces the word
by induction.
\end{proof}
The upper bound \eqref{eq:upper} follows by counting subsets of size at most $m$.
This is a time encoding, not a count of coordinate tuples in a radius-$n$ box.
It saves one power of $n$ per source cell relative to the crude planar count.

\begin{remark}[Exactly versus at most]
For fixed finite $n$, the sets of prefixes produced by exactly $m$ and by at
most $m$ initially black cells coincide: add padding outside the finite read
footprint. This is not an assertion about equality of their infinite trace
spaces. A padding cell may be reached later.
\end{remark}

\section{The complete infinite trace hierarchy}
\begin{lemma}[Bounded-cost branches have one finite source]\label{lem:branch}
Suppose each finite prefix of an infinite word is feasible and has minimal
source of size at most $m$. Then the word belongs to $\Tr_m$.
\end{lemma}
\begin{proof}
Minimal sources of nested prefixes are nested: an extension can add a first
black cell only outside the earlier read footprint. Their union has at most
$m$ elements. On every finite prefix, that union agrees with its minimal source
on all read cells. It therefore realizes every prefix and the entire stream.
\end{proof}
Thus $\Tr_m$ is closed in the compact product space. It is countable because
finite subsets of the countable lattice are countable. The infinite union in
Lemma~\ref{lem:branch} defines the unique minimal source $M(w)$ of a stream.
No algorithm deciding which cells of a given source are never read is assumed.

We recall the classical unboundedness argument, attributed to Bunimovich and
Troubetzkoy in the exposition of Gajardo, Lutfalla and Rao~\cite[Section~2]{highway}.
\begin{lemma}[Fresh sites occur arbitrarily late]\label{lem:fresh}
A classical ant orbit cannot visit only finitely many cells.
\end{lemma}
\begin{proof}
Otherwise let $I$ be the nonempty finite set visited infinitely often, and
choose $v\in I$ maximizing $x+y$. Its north and east neighbors eventually
are never visited. The parity of the heading axis plus checkerboard position
is invariant; hence every arrival at $v$ has a fixed axis. Late horizontal
arrivals must come from the west, heading east. Alternating reads at $v$
eventually force a north departure. Late vertical arrivals must come from
the south, heading north, and alternation forces an east departure. Either
case contradicts the absent neighbor. This is the usual corner argument.
\end{proof}
This lemma does not claim that individual cells are visited finitely often.

\begin{proof}[Proof of Theorem~\ref{thm:topology}]
If $|M(w)|=m$, choose a prefix that has read all its $m$ initially black cells.
Any stream in $\Tr_m$ with this prefix must contain exactly those required
source cells and can have no additional first black read. Its minimal source
is $M(w)$ and its entire future is fixed by determinism. Hence $w$ is isolated.

If $|M(w)|=r<m$, choose successive fresh cells of $w$, at times tending to
infinity and after all cells of $M(w)$ have been read. Blacken one of them in
addition to $M(w)$. The corresponding complete input agrees with $w$ until
that first visit, and disagrees at that visit. These distinct streams lie in
$\Tr_m$ and converge to $w$. Therefore the nonisolated points of $\Tr_m$ are
exactly $\Tr_{m-1}$. Iteration, with singleton $\Tr_0$, proves the theorem.
\end{proof}
In particular the point rank of $w\in\Tr_m$ is $m-|M(w)|$.
The prefix tree is decidable by same-cell alternation and first-read counting,
and every admitted finite node extends from its finite minimal source.
Compactness does not assert that this source is effectively identifiable from
an infinite stream after a uniformly bounded observation time.

\section{A whole-history protected binary channel}
This section supplies the ant-specific lower construction. The finite data
are printed in Appendix~\ref{app:data} and independently replayed. Let
\[
 t_0=11017,\quad p_0=(-35,-10),\quad d=(-2,-2),\quad q=(-2,-1),
 \qquad H=\{(x,y):x+y\le 0\}.
\]
Blank execution for $t_0$ updates has pose $(p_0,\text{west})$ and a complete black field
$B_*$ of size 835. Its part in $p_0+H$ is $p_0+C$, where $C$ is the displayed
21-cell set. It is the advancing front in the protected southwest half-plane:
the local identity below translates it by $d$ and leaves a finite residue
behind. Every other cell of $B_*$ is strictly outside that half-plane.
The blank prelude reads none of the ray $p_0+q+2 i d$, $i\ge 0$.

Let $w_{\rm start}$ denote the turn word of the first $t_0$ blank updates.
Let $F$ denote the one-update map on complete states $(B,p,h)$, and write
$0=(0,0)$ for the coordinate origin.
Let $W$ be the standard 104-letter word in the appendix. The same letter $W$
for a heading is always explicitly called ``west'' below. Let $S$ be its
228-letter single-cell perturbation word. Exact literal execution verifies
\begin{align}
 F^{104}(C,0,\text{west})
   &=((C+d)\sqcup R,d,\text{west}),\label{eq:kernel0}\\
 F^{228}(C\cup\{q\},0,\text{west})
   &=((C+2 d)\sqcup R_1,2 d,\text{west}).\label{eq:kernel1}
\end{align}
Every read in both equations lies in $H$. The entire residue $R$ has size 12
and lies in $-4<x+y\le 0$. The entire residue $R_1$ has size 25 and lies in
$-8<x+y\le 0$. In particular these are complete-field equations, not equality
modulo erased debris. Equation~\eqref{eq:kernel0} implies that $W$ repeats
forever, since each new copy stays in a strictly nested half-plane and never
reads an earlier residue.

Put $A_0=WW$ and $A_1=S$. Both blocks have the same final pose $2 d,\text{west}$
and the same retained front $C+2 d$. The zero-block residue is
$R_0=R\sqcup(R+d)$, of size 24; both residues are in $-8<x+y\le 0$.
Neither block reads any point $q+2 i d$ with $i\ge 1$. Both first read $q$ at
update 13, where their turns disagree; their preceding turns agree.

\begin{lemma}[Complete binary composition]\label{lem:binary}
For any finite $J\subset\N$, take the time-zero input
\begin{equation}\label{eq:seed}
 B_J=\{p_0+q+2 i d:i\in J\},
\end{equation}
with all other cells white. If $b_i=[i\in J]$, its entire turn stream is
\begin{equation}\label{eq:stream}
 w_{\rm start} A_{b_0} A_{b_1} A_{b_2}\cdots.
\end{equation}
After $L$ blocks its time and pose are
\[
 t_L=t_0+208L+20\sum_{i<L}b_i,\qquad
 p_L=p_0+2Ld,\qquad h_L=\text{west}.
\]
Its complete black field is
\begin{equation}\label{eq:full}
 \bigl(B_*\setminus(p_0+C)\bigr)
 \sqcup\bigsqcup_{i<L}(p_0+2 id+R_{b_i})
 \sqcup(p_L+C)
 \sqcup\{p_0+q+2 id:i\in J,\ i\ge L\}.
\end{equation}
Every such input eventually follows $W^\omega$.
\end{lemma}
\begin{proof}
All input markers avoid the startup footprint. At block $i$ the current
protected half-plane is $x+y\le p_{0,x}+p_{0,y}-8 i$. Every older residue lies
strictly outside it. The later markers lie inside but neither local block
reads them. Translate \eqref{eq:kernel0} twice or \eqref{eq:kernel1} according
to $b_i$. The complete residue falls into this block's own eight-level slab;
these slabs are disjoint. This proves \eqref{eq:full} and the pose/time
identities by induction. After the last marker only $WW$ blocks remain.
All markers were present at time zero; no intervention or replenishment occurs.
\end{proof}
For $|J|=m$, equation~\eqref{eq:full} has $835+24L+m$ black cells. The original
mass is $m$, whereas the evolving complete mass is unbounded. This illustrates
why a sparse-source complexity bound is not a bounded-memory assertion.

\begin{lemma}[Exact distinction time]\label{lem:distinction}
If two binary inputs first differ at slot $i$, with $r$ common earlier ones,
their turn streams first differ at update
\[
 t_0+208 i+20 r+13.
\]
\end{lemma}
\begin{proof}
All preceding blocks coincide. The two current blocks have a common prefix
of length 13 and opposite next turns, by the checked kernels.
\end{proof}

\section{Sharpness, topology inside the highway basin, and prediction}
\begin{proof}[Proof of Theorem~\ref{thm:prefix}]
The upper bound was proved in Section 2. For the lower bound choose all subsets
of $\{0,\ldots,L-1\}$ of size at most $m$ in \eqref{eq:seed}. Their first $L$
blocks finish by $t_0+208L+20\min(m,L)$. Lemma~\ref{lem:distinction} distinguishes
all the resulting prefixes by that time, and Lemma~\ref{lem:binary} proves each
input positive. This proves \eqref{eq:lower}. For fixed $m\ge 1$ take
$L=\max(0,\lfloor(n-t_0-20 m)/208\rfloor)$. The lower term $\binom Lm$ is
asymptotic to $n^m/(208^m m!)$ and the upper sum to $n^m/m!$. These prove the
sharp exponent, not equality of those constants.
\end{proof}

\begin{proof}[Proof of Theorem~\ref{thm:positive}]
Let $\mathcal S_m\subset\{0,1\}^{\N}$ contain all binary sequences with at
most $m$ ones. It is compact. Equations~\eqref{eq:seed}--\eqref{eq:stream}
define a continuous injection into the turn space: every finite observation
depends on finitely many slots. Lemma~\ref{lem:distinction} also gives a
continuous inverse on the image, which we call $\Pos_m$. The image is compact
and consists entirely of positive inputs. In $\mathcal S_m$, a point with
exactly $m$ ones is isolated; every point with fewer ones is approached by
adding a one at a coordinate tending to infinity. Thus
$\mathcal S_m'=\mathcal S_{m-1}$, giving the same identity for $\Pos_m$.
The all-zero source produces the blank orbit, so $\Pos_0=\Tr_0$.
The prefix lower bound is the one just proved.
\end{proof}

For prediction we count incorrect next-turn guesses, using the usual
worst-case mistake criterion~\cite{littlestone}.
\begin{corollary}[Sharp passive mistake budget]\label{cor:prediction}
With unknown initial source and observed past turns, the optimal worst-case
number of mistakes of a deterministic next-turn predictor on $\Tr_m$ is $m$.
This remains a lower bound when inputs are restricted to $\Pos_m$.
\end{corollary}
\begin{proof}
The observed turns determine the current position and all earlier reads.
Predict the forced color on revisits and $\Rr$ on first visits. Exactly the
first reads of initially black cells cause mistakes, at most $m$ of them.
For any competing deterministic predictor, choose $m$ successive binary
slots. At the decisive update 13 in each slot, choose the source bit opposite
to its prediction. The decisions are well-defined from the preceding chosen
history. They specify one complete finite seed with at most $m$ black cells,
not interventions during execution. Lemma~\ref{lem:binary} realizes it and
Lemma~\ref{lem:distinction} supplies $m$ mistakes.
\end{proof}
The predictor can store an unbounded visited-cell map. This statement neither
supplies a finite-state predictor nor accelerates a known input's simulation.
For a supplied $n$, distinguishing all prefixes needs
$m\log_2 n+O_m(1)$ bits of advice: the upper bound encodes first-black times,
and the positive lower family gives the matching order.

\section{Limitations}
The upper bound and derivative identity apply to every complete finite source
in the stated mass class, regardless of which HC branch it might follow.
They do not assume convergence, classify possible periodic wakes, or require a
strip. The protected family is only the witness of sharpness, not a substitute
for that all-input scope.

The result is not a termination rank. Shifting the blank stream by four updates
gives a word beginning $\Ll$, whereas the unique member of $\Tr_0$ begins
$\Rr$. Thus $\Tr_0$ is not shift-invariant; more generally an update can increase
the mass of the complete current state. The derived-set rank belongs to a
space of inputs at a specified origin of time. It does not decrease along a
physical orbit. Zero exponential anchored prefix rate and finite rank therefore
do not imply eventual periodicity.

The theorem does not assert that the full positive subset of $\Tr_m$ is closed;
only the explicitly constructed $\Pos_m$ is compact. No existence or exclusion
of a finite HC counterexample is obtained. The open all-input question still
asks whether a single finite source can support infinitely many nonstandard
events. The present census answers how many observable histories that family
can contain, not whether every one of them standardizes.

\section{Reproduction and evidence}
Run \texttt{python sparse\_trace\_space/reproduce.py} from the repository
root. The default copies only the theorem-local package into a temporary
directory, recomputes a deterministic report, and compares it with the frozen
one. It does not rewrite evidence. Python standard library and C++17 suffice.
The kernels, source fields and complete residues are explicit in the supplement;
its source hashes identify the actual checked inputs. They are available in
the \href{https://github.com/kehao95/langtons-ant-core/tree/3ae4ce12abf6ad44a5c7625b31d25e6f1365ffa3/sparse_trace_space}{immutable public sparse-trace source package}.

The verification reconstructs the complete first-read tree to depth 24 with
source budget three and checks every time-code encoding and decoding.
Independent heading-index Python and vector-turn C++ simulations compare
complete endpoints for every binary source word of length at most eight,
every field on a $3\times3$ set, and additional sources with remote unread
inventory. Every pair of eight-bit sources has the prescribed first distinction
time. Further controls check fresh-site approximation, exact integer counting,
huge symbolic slots and malformed inputs. The source-bound report records the
individual cases and counts. Infinite and all-$m$ claims follow from the
coverage and geometric induction above.

\appendix
\section{Complete local data}\label{app:data}
All listed coordinates are relative to a west-facing head at the origin.
The 21-cell front is
\[
\begin{split}
C=\{&(-2,-4),(-2,-3),(-2,0),(-2,1),(-1,-5),(-1,-2),(-1,-1),\\
&(-1,0),(-1,1),(0,-6),(0,-5),(0,-4),(0,0),(1,-6),(1,-5),\\
&(1,-4),(1,-3),(1,-2),(2,-5),(3,-4),(3,-3)\}.
\end{split}
\]
The standard word, with line breaks omitted, is
\begin{quote}\ttfamily\small
LRLLRRRRLLRLLRRRRLLRRRRLLRLRRRRLRLLLLRRRRLRRLRRRRLLL\\
LRLRRRRLRRRRLLLLRLRRRRLRLLRRLLLLRRLLRRRRLLRRLRLLRLLR
\end{quote}
The perturbed word $S$, with line breaks omitted, is
\begin{quote}\ttfamily\small
LRLLRRRRLLRLLLLRRRRLRLRRRRLRRRRLLLLRLRRRRLRRRRLLLLRL\\
RLLLLRRRRLRRLRRRRLLLLRLRRRRLRRRRLLLLRLRRRRLRLLRRLLLL\\
RRLLRRRRLLRRLRLLRLLRLRLLRRRRLLRLLRRRRLLRRRRLLRLRRRRL\\
RLLLLRRRRLRRLRRRRLLLLRLRRRRLRRRRLLLLRLRRRRLRLLRRLLLL\\
RRLLRRRRLLRRLRLLRLLR
\end{quote}
The complete residue after $W$ is
\[
\begin{split}
R=\{&(-3,0),(-2,-1),(-2,1),(-1,-2),(-1,1),(0,-3),\\
 &(0,-2),(0,0),(1,-4),(2,-5),(3,-4),(3,-3)\}.
\end{split}
\]
The complete residue after $S$ is
\[
\begin{split}
R_1=\{&(-5,-2),(-5,-1),(-5,0),(-4,-3),(-4,1),(-3,-4),(-3,-1),\\
&(-3,2),(-2,-5),(-2,-4),(-2,-2),(-2,2),(-1,-6),(-1,-2),\\
&(-1,1),(0,-7),(0,-3),(0,-2),(0,0),(1,-6),(1,-5),(1,-4),\\
&(2,-5),(3,-4),(3,-3)\}.
\end{split}
\]
These tuples and words determine the full finite checks
\eqref{eq:kernel0}--\eqref{eq:kernel1}. The 835-cell startup state is generated
by exactly 11017 literal updates from the empty source, rather than specified
by a truncated window. Each startup read is tested against the entire marker
ray: membership requires equal coordinate differences, divisibility by 4,
and a nonnegative index. The two blocks are checked the same way.

\begin{thebibliography}{9}
\bibitem{gajardo} A.~Gajardo. A symbolic projection of Langton's Ant
(Extended Abstract). \emph{DMTCS Proceedings AB, DMCS'03}, 57--68, 2003.
\href{https://doi.org/10.46298/dmtcs.2312}{doi:10.46298/dmtcs.2312}.
\bibitem{oneblack} H.~Ke. \emph{Every One-Black-Cell Langton Ant Reaches the
Period-104 Highway: A Lean-Checked Proof}. 2026.
\href{https://github.com/kehao95/langtons-ant-core/tree/08fd2a5fd07dde19ecf0177498e07463f2bb55f9/one_black}{Article and proof project}.
The phase-13, 228-step channel is also recorded in the original single-obstacle
and separated-obstacle certificate families; the local data required here are
fully reproduced in Appendix~\ref{app:data}.
\bibitem{salo} V.~Salo and I.~T\"orm\"a.
\emph{Complexity of Conjugacy, Factoring and Embedding for Countable Sofic
Shifts of Rank 2}. Author preprint, 2014.
\href{https://arxiv.org/abs/1408.6695v1}{arXiv:1408.6695v1}.
\bibitem{littlestone} N.~Littlestone.
Learning quickly when irrelevant attributes abound: A new linear-threshold algorithm.
\emph{Machine Learning} 2 (1988), 285--318.
\href{https://doi.org/10.1007/BF00116827}{doi:10.1007/BF00116827}.
\bibitem{highway} A.~Gajardo, V.~H.~Lutfalla and M.~Rao.
Ants on the highway. \emph{Natural Computing} 24 (2025), 497--509.
\href{https://doi.org/10.1007/s11047-025-10018-9}{doi:10.1007/s11047-025-10018-9}.
Author version: \href{https://arxiv.org/abs/2409.10124v1}{arXiv:2409.10124v1}.
\end{thebibliography}
\end{document}
